% Desarrollo acotado: lector energético intervalar del registro K.
% Estatuto: formalización añadida sobre K y sobre las identidades de Schur
% ya presentes en el corpus. No se atribuye novedad histórica a Schur.
\section{Formas de Dirichlet y equivalencia variacional del refinamiento}
\label{sec:energia-refinamiento-k}

La representación positiva del registro dodecafásico define una partición
ordenada mediante
\[
 K=(234,543,140,729,659,824,621,58,914,794,146,601),
 \qquad d_j=K_j/6263,
 \qquad t_i=\sum_{j=1}^id_j.
\]
El registro se recibe de la composición
\(x_{\rm term}\mapsto\mathcal R_{12}\mapsto
(b^{90},b^{120},Q_{\rm TPK})\mapsto U_{\rm sgn}\mapsto K\),
posterior a APP--TRIT--TPK y a su estado enriquecido, conforme a los lectores
reconstruidos en \cite{hmtX}. Esta sección aplica
un lector cuadrático a la partición resultante. La comparación entre
profundidades se efectuará por restricciones a extremos heredados, de modo
que cada coeficiente de la energía conserve su origen en una separación
positiva ya producida.

El corpus contiene una energía ponderada de diferencias para el grafo de
las dos memorias y una reducción de Schur de las formas completas. Aquí se
reúnen esas operaciones sobre la red intervalar del registro. La
especialización a conductancias \(1/d_j\), el operador de interpolación
y su compatibilidad con la deformación de las separaciones constituyen
la formalización desarrollada en esta sección.

\subsection{Reconstrucción de las separaciones a partir de la matriz de rigidez}

Sea \(d=(d_1,\ldots,d_m)\), con \(d_j>0\), y sean
\(f=(f_0,\ldots,f_m)\in\mathbb C^{m+1}\) los valores de un observable
en los extremos de la partición. Definimos
\begin{equation}
 \mathcal E_d(f)=\sum_{j=1}^{m}\frac{|f_j-f_{j-1}|^2}{d_j},
 \qquad L_d=D^*\operatorname{diag}(d_1^{-1},\ldots,d_m^{-1})D,
 \label{eq:energia-intervalar}
\end{equation}
donde \(D:\mathbb C^{m+1}\to\mathbb C^m\) es el operador de
incidencia orientada, \((Df)_j=f_j-f_{j-1}\), y el asterisco denota la
adjunción hermítica. Así \(\mathcal E_d(f)=f^*L_df\).
La matriz de rigidez tiene entradas
\[
 (L_d)_{00}=d_1^{-1},\quad (L_d)_{mm}=d_m^{-1},\quad
 (L_d)_{ii}=d_i^{-1}+d_{i+1}^{-1}\quad(1\le i<m),
\]
\[
 (L_d)_{i-1,i}=(L_d)_{i,i-1}=-d_i^{-1},
\]
y las demás son nulas. Es hermítica positiva semidefinida, su núcleo
consiste en las funciones constantes y su forma es positiva definida
en el cociente por ese núcleo. La matriz completa recupera las
separaciones mediante
\begin{equation}
 d_j=-\frac1{(L_d)_{j-1,j}}.
 \label{eq:recuperacion-d-rigidez}
\end{equation}

\begin{proposition}[Equivalencia entre separaciones y lector nodal completo]
\label{prop:energia-bidireccional}
La aplicación \(d\mapsto L_d\) es una biyección entre
\((0,\infty)^m\) y el conjunto de matrices reales simétricas
tridiagonales cuyas entradas inmediatamente adyacentes a la diagonal son
negativas y cuya suma en cada fila es cero. Su inversa es
\eqref{eq:recuperacion-d-rigidez}. La normalización \(\sum_jd_j=1\)
equivale a \(\sum_{j=1}^m-1/(L_d)_{j-1,j}=1\).
\end{proposition}
\begin{proof}
Las fórmulas explícitas de \(L_d\) verifican todas las condiciones.
Recíprocamente, las entradas adyacentes negativas determinan separaciones
positivas únicas por \eqref{eq:recuperacion-d-rigidez}. La suma nula de
cada fila determina sus entradas diagonales como la suma de las
conductancias incidentes. La tridiagonalidad fija en cero las restantes
entradas. Se recupera exactamente la matriz de
\eqref{eq:energia-intervalar}; su semidefinitud positiva sigue entonces
de la factorización \(D^*\operatorname{diag}(1/d_j)D\).
\end{proof}

La derivada discreta ponderada
\(J_j=(f_j-f_{j-1})/d_j\) es el flujo orientado de la arista \(j\).
En los vértices interiores, \((L_df)_i=J_i-J_{i+1}\); en los extremos,
\((L_df)_0=-J_1\) y \((L_df)_m=J_m\). La ecuación interior
\(L_df=0\) expresa que el flujo tiene el mismo valor en cada arista.
Esta definición es la de un observable cuadrático del lector intervalar;
su realización como magnitud física requiere la coordenatización dimensional pertinente.

\subsection{Interpolación armónica, complemento de Schur y reconstrucción del detalle}

Refinemos cada separación como
\(d_j=\sum_{c=1}^{r_j}\delta_{j,c}\), con \(\delta_{j,c}>0\).
Sea \(\widetilde d\) la concatenación de esas separaciones, \(R\) la
restricción a los extremos antiguos y \(H\) la interpolación definida por
\begin{equation}
 (Hf)_{j,c}=(1-\theta_{j,c})f_{j-1}+\theta_{j,c}f_j,
 \qquad
 \theta_{j,c}=\frac{\sum_{a=1}^c\delta_{j,a}}{d_j}.
 \label{eq:interpolacion-armonica-k}
\end{equation}
Los extremos compartidos se cuentan una sola vez. Se cumple \(RH=I\).

\begin{theorem}[Refinamiento isométrico de la energía]
\label{thm:energia-schur-refinamiento}
Para todo vector de valores antiguos \(f\),
\begin{equation}
 H^*L_{\widetilde d}H=L_d,
 \qquad L_{\widetilde d}H=R^*L_d,
 \qquad
 \min_{Rg=f}\mathcal E_{\widetilde d}(g)=\mathcal E_d(f).
 \label{eq:isometria-energia-k}
\end{equation}
El minimizador es único y vale \(Hf\). Todo \(g\) admite la
descomposición única
\begin{equation}
 g=H(Rg)+z,\qquad Rz=0,
 \qquad
 \mathcal E_{\widetilde d}(g)
 =\mathcal E_d(Rg)+\mathcal E_{\widetilde d}(z).
 \label{eq:detalle-energia-k}
\end{equation}
\end{theorem}
\begin{proof}
En una arista antigua escribamos \(h_c=g_{j,c}-g_{j,c-1}\), de modo
que \(\sum_ch_c=f_j-f_{j-1}=a\). Si
\(h_c=\delta_{j,c}a/d_j+r_c\), entonces \(\sum_cr_c=0\) y
\[
 \sum_c\frac{|h_c|^2}{\delta_{j,c}}
 =\frac{|a|^2}{d_j}+\sum_c\frac{|r_c|^2}{\delta_{j,c}}.
\]
La igualdad se obtiene expandiendo; el término cruzado es proporcional a
\(\sum_cr_c\) y se anula. El mínimo se alcanza exactamente cuando
todos los \(r_c\) son cero, condición que determina
\eqref{eq:interpolacion-armonica-k}. Al sumar las aristas se obtienen
el mínimo y la descomposición energética. El flujo de \(Hf\) es
constante dentro de cada arista antigua, por lo que sus componentes
interiores de \(L_{\widetilde d}Hf\) son cero y sus componentes heredadas
coinciden con \(L_df\). Esto prueba la segunda identidad matricial; al
multiplicar por \(H^*\), usando \(RH=I\), se obtiene la primera.
\end{proof}

El término \(\mathcal E_{\widetilde d}(z)\) es no negativo y se anula
exactamente para \(z=0\): un detalle de energía nula sería constante y,
al anularse en los extremos heredados, se anularía en toda la red.
La pareja \((Rg,z)\) conserva el vector refinado completo y permite
recuperarlo por \(g=H(Rg)+z\). La minimización selecciona su
componente armónica; el archivo del detalle \(z\) mantiene la
información que la sola forma reducida no distingue. La conservación
de la energía observada y la conservación del estado refinado quedan
así expresadas por mapas diferentes y explícitos.

Ordenando los vértices como extremos heredados \(E\) e interiores \(I\),
escribimos
\[
 L_{\widetilde d}=\begin{pmatrix}A&B^*\\B&C\end{pmatrix}.
\]
Para un refinamiento con vértices interiores, el bloque \(C\) es
positivo definido: una función nula en los extremos
que tenga energía cero es constante en cada subintervalo y, por tanto,
idénticamente nula. Se obtiene
\begin{equation}
 H=\begin{pmatrix}I\\-C^{-1}B\end{pmatrix},\qquad
 \operatorname{Schur}_{I}(L_{\widetilde d})
 =A-B^*C^{-1}B=L_d.
 \label{eq:schur-intervalar-k}
\end{equation}
Si el refinamiento conserva todos los vértices y no introduce ninguno,
la misma afirmación se reduce a \(H=I\) y \(L_{\widetilde d}=L_d\).

\begin{proposition}[Naturalidad de las prolongaciones y de la eliminación]
\label{prop:energia-naturalidad-sucesiva}
Para tres particiones anidadas, las interpolaciones armónicas y los
complementos de Schur satisfacen
\begin{equation}
 H_{n+2\leftarrow n}
 =H_{n+2\leftarrow n+1}H_{n+1\leftarrow n},
 \qquad
 \operatorname{Schur}_{n+1\to n}
 \bigl(\operatorname{Schur}_{n+2\to n+1}(L_{n+2})\bigr)=L_n.
 \label{eq:energia-naturalidad-sucesiva}
\end{equation}
\end{proposition}
\begin{proof}
La interpolación de valores afines sobre un intervalo sigue siendo la
misma función afín después de cualquier subdivisión de ese intervalo.
Esto prueba la primera identidad. Minimizar primero sobre los vértices
incorporados en el nivel último y después sobre los del nivel intermedio
es minimizar sobre todos ellos, con los mismos extremos heredados fijados.
La identidad variacional de \eqref{eq:isometria-energia-k} y la unicidad
del minimizador prueban la segunda identidad. La restricción del mínimo
se realiza, por tanto, mediante la misma composición que la prolongación
armónica de los datos.
\end{proof}

\subsection{Respuesta de Dirichlet a Neumann y defecto energético de una extensión}

Si sólo se conservan los extremos inicial y final, sea
\(D_0=\sum_{j=1}^md_j\). Para valores \((a,b)\), la solución
armónica es
\[
 f_i=a+(b-a)\frac{\sum_{j=1}^id_j}{D_0},\qquad
 J_j=(b-a)/D_0.
\]
Su operador de Dirichlet a Neumann, que transforma valores de frontera en
flujos salientes, es
\begin{equation}
 \Lambda_d=\frac1{D_0}
 \begin{pmatrix}1&-1\\-1&1\end{pmatrix},\qquad
 \min\mathcal E_d=\frac{|b-a|^2}{D_0}.
 \label{eq:dtn-intervalar-k}
\end{equation}
Cuando se prescribe el valor en todos los vértices heredados, su operador
de respuesta multipuntual es, por el mismo argumento, \(\Lambda=L_d\).

Una extensión aproximada \(g=(f,y)\), donde \(y\) reúne los valores
en los vértices interiores, tiene residual interior
\(r=Bf+Cy\). La identidad exacta
\begin{equation}
 \mathcal E_{\widetilde d}(g)-f^*\Lambda f
 =r^*C^{-1}r\ge0
 \label{eq:residual-dtn-k}
\end{equation}
se deduce de \(y+C^{-1}Bf=C^{-1}r\) y de la descomposición
\[
 g^*L_{\widetilde d}g
 =f^*(A-B^*C^{-1}B)f
 +(y+C^{-1}Bf)^*C(y+C^{-1}Bf).
\]
Así, el valor efectivo y el detalle mantienen una descomposición
hermítica completa. En particular, si una
extensión lineal es \(\widetilde H=\begin{pmatrix}I\\Z\end{pmatrix}\), su matriz de
residuales es \(\mathscr R=B+CZ\) y
\[
 \widetilde H^*L_{\widetilde d}\widetilde H-\Lambda
 =\mathscr R^*C^{-1}\mathscr R\succeq0.
\]
El residual nulo caracteriza la interpolación armónica exacta. Este
residual energético vive en el espacio de vértices interiores; la
cancelación de residuos de formas meromorfas pertenece a otro mapa de
frontera y conserva su definición propia.

\subsection{Conmutación del refinamiento y del flujo de separaciones}

Sea \(u=P_3P_{11}K\) la dirección algebraica ya recuperada de los
lectores del registro \cite{hmtX}. Para el parámetro real \(s\), definamos
\[
 d_j(s)=\frac{K_je^{su_j}}{\sum_\ell K_\ell e^{su_\ell}}.
\]
Fijadas fracciones positivas \(\vartheta_{j,c}\) de suma uno para cada
\(j\), el refinamiento heredado es
\[
 \delta_{j,c}(s)=\vartheta_{j,c}d_j(s),\qquad
 u_{j,c}=u_j.
\]
Producir primero los pesos refinados \(K_j\vartheta_{j,c}\) y
deformarlos da la misma colección: los sucesores tienen idéntico exponencial y
sus coeficientes suman \(K_j\). Para todo \(s\),
\begin{equation}
 \operatorname{Schur}_{I}(L_{\widetilde d(s)})=L_{d(s)},
 \qquad H^*L_{\widetilde d(s)}H=L_{d(s)}.
 \label{eq:schur-flujo-k}
\end{equation}
En esta comparación padre--hijos, \(H\) es constante: sus coeficientes
son las sumas parciales de \(\vartheta_{j,c}\). En cambio, la
interpolación desde los dos extremos globales usa
\(t_i(s)=\sum_{j\le i}d_j(s)\) y cambia con el parámetro. Ambas
afirmaciones corresponden a restricciones diferentes de la misma red.

La normalización \(\sum_jd_j(s)=1\) hace que el operador global de dos
extremos permanezca igual a
\(\left(\begin{smallmatrix}1&-1\\-1&1\end{smallmatrix}\right)\).
La matriz de todos los nodos \(L_{d(s)}\) sí conserva la deformación y
la recupera mediante \eqref{eq:recuperacion-d-rigidez}. La pérdida de
información en el lector de dos extremos queda identificada exactamente:
retiene la longitud total y sus flujos, mientras el lector nodal retiene
todas las separaciones. El archivo de detalles proporciona el complemento
necesario para reconstruir el observable refinado.
